\documentclass[a4paper, 12pt]{report}

\usepackage[T2A]{fontenc}
\usepackage[utf8]{inputenc}
\usepackage[english, russian]{babel}
\usepackage{enumitem}
\usepackage{xcolor}
\usepackage{listings}
\setlist[itemize]{label=\textcolor{black}{\textbf{\textendash}}, itemsep=10pt}

\author{Чахкиев Руслан}
\title{Тех.Документация BrickGame v1.0 aka Tetris}
\date{9 ноября 2025 г.}

\begin{document}

\maketitle
\tableofcontents{}

\clearpage
% =====НОВАЯ ГЛАВА=====
\chapter{Общие сведения}

\begin{itemize}
    \item \textbf{Название проекта:} BrickGame v1.0
    \item \textbf{Назначение:} Консольная игра «Тетрис» с управлением с помощью клавиатуры
    \item \textbf{Платформа:} Linux (через терминал)
    \item \textbf{Разработчик:} Чахкиев Руслан
    \item \textbf{Версия:} 1.0
\end{itemize}

\clearpage
% ======НОВАЯ ГЛАВА=====
\chapter{Цель и Задачи}

\section{Цель}
Реализовать классическую игру "Тетрис" в консоли.

\section{Основные задачи}
\begin{itemize}[label=\textcolor{black}{\textbf{\textendash}}]
    \item Реализовать падение фигур с определённой периодичностью
    \item Движение фигуры вправо, влево, вниз во время игры
    \item Вращение фигуры на 90 градусов по часовой стрелке
    \item Подсчет уровня и очков
    \item Система рекорда, который будет храниться в отдельном файле
    \item Отображение текущего состояния игры
\end{itemize}

\clearpage
% =====НОВАЯ ГЛАВА=====
\chapter {Структура программы}

\section{Архитектура программы}
Программа реализована в соответствии с паттерном \textbf{MVC (Model View Controller)}.
То есть, программа разделенна на 
\begin{itemize}
    \item \textbf{Бекэнд(Логика игры):} brick\_game/tetris
    \item \textbf{Контроллер(Ввод клавиши и.т.п):} brick\_game/tetris/controller
    \item \textbf{Фронтенд/(Отрисовка):} gui/cli
\end{itemize}

А также в программе присутствует система \textbf{КА (Конечных Автоматов)}.\\
То есть, в зависимости от введённой пользователем клавиши, устанавливается определённое действие.\\
В зависимости от действия, устанавливается определённое состояние\\

\clearpage
\section{Структура программы}

\begin{itemize}
    \item \textbf{brick\_game}
    \begin{itemize}
        \item \textbf{tetris}
        \begin{itemize}
            \item \textbf{controller}
            \begin{itemize}
                \item s21\_tetris\_controller.c
                \item s21\_tetris\_controller.h
            \end{itemize}
            \item s21\_back\_field\_actions.c
            \item s21\_back\_figure\_actions.c
            \item s21\_back\_game\_info.c
            \item s21\_back\_init.c
            \item s21\_back\_matrix\_actions.c
            \item s21\_backend.h
        \end{itemize}
    \end{itemize}
    
    \item \textbf{gui}
    \begin{itemize}
        \item \textbf{cli}
        \begin{itemize}
            \item s21\_front\_init\_funcs.c
            \item s21\_front\_print\_actions.c
            \item s21\_frontend.c
            \item s21\_frontend.h
        \end{itemize}
    \end{itemize}
\end{itemize}

\begin{verbatim}
Взаимодействие:
Главный цикл -> Обновление логики -> Перерисовка экрана -> Ввод
Цикл         -> Бекэнд            -> Фронтенд           -> Контроллер
\end{verbatim}

\clearpage
% =====НОВАЯ ГЛАВА=====
\chapter{Основные алгоритмы}

\begin{itemize}
    \item Падение фигуры: Каждые N миллисекунд сдвиг вниз.
    \item Поворот фигуры на 90 градусов.
    \item Проверка столкновений: Если фигура достигла дна или натыкается на следующую фигуру, то генерируется новая фигура.
    \item Удаление линий: Если строка полностью заполнена, строка очищается, и все линии поля сдвигаются вниз на 1 клетку.
    \item Подсчёт очков: 100 за одну удалённую линию, 300 - за две, 700 - за три, 1500 - за четыре.
    \item Генерация фигур: случайным образом из 7 типов.
    \item Отображение и сохранение рекордов в файле \textbf{record.bin}.
\end{itemize}

\clearpage
% =====НОВАЯ ГЛАВА=====
\chapter{Основные структуры данных}

\section{Основная структура \texttt{GameData\_t}}

Хранит в себе \textbf{все} данные об игре
\begin{verbatim}
    typedef struct {
        int **figure[FIGURE_COUNT];
        Figure_t fig_cur;
        Figure_t fig_next;
        GameInfo_t g_i;
        GameState_t g_s;
        UserAction_t g_a;
    } GameData_t;
\end{verbatim}

\section{Структура \texttt{GameInfo\_t}}

Хранит в себе параметры об игре (Текущий счет, рекорд и.т.п)
\begin{verbatim}
    typedef struct {
        int **field;
        int **next;
        int score;
        int high_score;
        int level;
        int speed;
        int pause;
    } GameInfo_t;
\end{verbatim}

\section{Структура \texttt{FIgure\_t}}

Хранит в себе информацию о текущей и следующей фигуре
\begin{verbatim}
    typedef struct {
        int **shape;
        int type;
        f_cords cord;
    } Figure_t;

    typedef struct {
        int y, x;
    } f_cords;
\end{verbatim}

\section{Структура \texttt{GameAction\_t}}

Хранит в себе информацию о текущем действии
\begin{verbatim}
    typedef enum {
        Start,
        Pause,
        Terminate,
        Left,
        Right,
        Up,
        Down,
        Action,
        Drop,
        Nothing
    } UserAction_t;
\end{verbatim}

\section{Структура \texttt{GameState\_t}}

Хранит в себе информацию о текущем состоянии
\begin{verbatim}
    typedef enum {
        STATE_MENU,
        STATE_PLAYING,
        STATE_PAUSE,
        STATE_GAME_OVER,
        STATE_EXIT
    } GameState_t;
\end{verbatim}

\clearpage
% =====НОВАЯ ГЛАВА=====
\chapter{Интерфейс}

\begin{itemize}
    \item Управление с клавиатуры.
    \begin{itemize}
        \item ← / → — Перемещение фигуры
        \item ↓     — Ускорение
        \item ↑     — Поворот
        \item R     — Начать игру
        \item Q     — Выход
        \item P     — Пауза
        \item SPACE — Дроп
    \end{itemize}
    \item Игровое поле отображается в терминале.
    \item Очки и уровень выводятся сбоку.
\end{itemize}

\clearpage
% =====НОВАЯ ГЛАВА=====
\chapter{Инструкция по запуску}

\section{Инструменты, необходимые для сборки}
Требуются две \textbf{сторонние} библиотеки
\begin{itemize}
    \item \textbf{ncurses.h} (Библиотека, для работы с графикой в терминале)
    \item \textbf{check.h} (Библиотека, для запуска тестирования)
    \item \textbf{Tex Live} (Дистрибутив для компиляции из \textbf{.tex} в \textbf{.pdf})
\end{itemize}

\section{Скачивание инструментов, необходимых для сборки}
Чтобы их скачать, в терминале нужно ввести\\
sudo apt install check \textbf{(Для check.h)}\\
sudo apt install libncurses5-dev libncursesw5-dev \textbf{(Для ncurses.h)}\\
sudo apt install texlive-full \textbf{(Для Tex Live)}\\
Также для компиляции из \textbf{.tex} в \textbf{.pdf} можно использовать онлайн-сервисы, например \textbf{"Overleaf"}\\
"https://ru.overleaf.com/project"

\clearpage
% =====НОВАЯ ГЛАВА=====
\chapter{Цели в Makefile}

\begin{itemize}
    \item \textbf{all}
        Собирает исполняемый файл для игры в Тетрис.
    \item \textbf{test}
        Собирает исполняемый файл для запуска тестов.
    \item \textbf{dvi}
        Компилирует документацию из \textbf{.tex} в \textbf{.pdf}.
    \item \textbf{dist}
        Создаёт архив с файлами, необходимыми для сборки исполняемого файла для игры в Тетрис.
    \item \textbf{gcov\_report}
        Генерирует html страницу, в которой указано сколько \textbf{строк} и \textbf{функций} были покрыты тестами.
    \item \textbf{install}
        Копирует исполняемый файл для игры в Тетрис в папку \textbf{BrickGame v1.0 aka Tetris}.
    \item \textbf{uninstall}
        Удаляет исполняемый файл для игры в Тетрис из папки \textbf{BrickGame v1.0 aka Tetris}.
    \item \textbf{clang\_check} и \textbf{clang\_fix}
        Проверяет соответствие написание кода с GoogleStyle и соответственно фиксит несоответствие.
    \item \textbf{leak\_game} и \textbf{leak\_test}
        Проверяет на утечку памяти игру и тесты.
    \item \textbf{clean}
        Очищает всё, кроме файлов, необходимых для сборки проекта.
\end{itemize}

\clearpage
% =====НОВАЯ ГЛАВА=====
\chapter{Заключение}

\begin{verbatim}
Проект реализует основные механики классического Тетриса.
Программа работает стабильно, код модульный и легко расширяемый.
\end{verbatim}

\end{document}